% !PRINTED 263
\weberpage{263}

If we now consider the sequence of quantities from this point of view,
\[
R_m,\ R_{m-1},\ R_{m-2} \ldots R_1,\ R_0,
\]
then in passing from $R_k$ to $R_{k-1}$ there occurs a permanence of sign or a change of sign, according as the quotient $R_k : R_{k-1}$ is positive or negative.

We can therefore also state the following theorem:

V.\ \gesperrt{If $\pi$ is the number of positive, $\nu$ that of the negative squares of $\varphi$, then $\pi$ is equal to the number of permanences of sign, $\nu$ equal to the number of changes of sign in the chain}
\weq{4}{R_m,\ R_{m-1},\ R_{m-2} \ldots R_1,\ R_0.}

This formulation of the theorem has the advantage that it can be carried over to the case in which individual interior terms of the sequence (4) vanish, provided only that no two consecutive terms are zero (which by \S.\,82, III. may always be assumed). For if $R_k = 0$, while $R_{k-1}$, $R_{k+1}$ are different from zero, then by \S.\,82, Theorem IV, $R_{k-1}$ and $R_{k+1}$ have different signs. In the sequence
\[
R_{k+1},\ R_k,\ R_{k-1}
\]
there thus occurs one change of sign and one permanence of sign, no matter whether we replace the vanishing $R_k$ by a positive or by a negative quantity.

If now we replace the chain (4)
\[
R_m,\ R_{m-1} \ldots R_1,\ R_0,
\]
in which individual terms vanish, by another,
\weq{5}{R'_m,\ R'_{m-1} \ldots R'_1,\ R'_0,}
in which no term vanishes, and in which to the non-vanishing terms of the sequence (4) there correspond terms of the same sign, then the sequences (4) and (5) have equally many changes of sign, whatever signs the $R'$ corresponding to the vanishing $R$ may have.

Now let $\psi(x)$ be a second arbitrary quadratic form of the variables $x$, with which we form the form
\weq{6}{\varphi' = \varphi + \varepsilon \psi}
in which $\varepsilon$ is a still undetermined coefficient. We shall now show at once that we can choose $\varepsilon$ in such a way that $\varphi'$ and $\varphi$ have the same number of positive and negative squares
%CONT>>
